% BHU PYQ -- Manually transcribed & error-corrected
% Paper: GLB-101: Elementary Physical and Structural Geology
% Exam: B.Sc. (Hons.) Semester I Examination, 2022-23

\documentclass[11pt,a4paper]{article}
\usepackage[a4paper,top=2.5cm,bottom=2.5cm,left=2.8cm,right=2.8cm,headheight=14pt]{geometry}
\usepackage{amsmath,amssymb}
\usepackage[T1]{fontenc}
\usepackage[utf8]{inputenc}
\usepackage{lmodern,microtype,fancyhdr,enumitem,hyperref,lastpage,parskip}

\hypersetup{colorlinks=true,urlcolor=black,linkcolor=black,
  pdftitle={GLB-101 Elementary Physical and Structural Geology -- BHU B.Sc. Sem I 2022-23}}
\pagestyle{fancy}\fancyhf{}
\renewcommand{\headrulewidth}{0.4pt}\renewcommand{\footrulewidth}{0.4pt}
\fancyhead[L]{\small\textit{Banaras Hindu University}}
\fancyhead[R]{\small\textit{GLB-101: Elementary Physical \& Structural Geology}}
\fancyfoot[C]{\small Page \thepage\ of \pageref{LastPage}}

\newlist{parts}{enumerate}{2}
\setlist[parts,1]{label=(\alph*),leftmargin=*,itemsep=5pt,topsep=3pt}
\newcommand{\pts}[1]{\hfill\textit{[#1]}}

\begin{document}
\begin{center}
  {\large\bfseries BANARAS HINDU UNIVERSITY}\\[2pt]
  {\normalsize B.Sc. (Hons.) --- Semester I Examination, 2022-23}\\[6pt]
  \rule{0.8\linewidth}{0.5pt}\\[6pt]
  {\large\bfseries Paper: GLB-101 --- Elementary Physical and Structural Geology}\\[4pt]
  {\normalsize Time: 3 Hours \hfill Maximum Marks: 70}\\[2pt]
  {\small Ref. No.: 055637}
\end{center}
\rule{\linewidth}{0.4pt}
\smallskip
\noindent\textit{Instructions: Answer \textbf{five} questions, selecting at least \textbf{two} each from Sections A and B, and Section-C which is compulsory.}
\bigskip

\bigskip\noindent{\large\bfseries SECTION --- A}
\rule{\linewidth}{0.4pt}\smallskip

\noindent\textbf{Question 1.} With a neat block diagram of a fault, define different elements of it. Describe different types of fault based on the rake of the net slip and attitude of fault relative to the attitude of the adjacent rocks with neat diagrams. \pts{14}

\medskip\noindent\textbf{Question 2.} Describe the difference between the conformable strata and unconformity in a sequence of strata. Discuss the different types of unconformities along with their origin with neat diagrams. \pts{14}

\medskip\noindent\textbf{Question 3.} Discuss how physical processes control weathering and erosion of a rock. \pts{14}

\medskip\noindent\textbf{Question 4.} Discuss the discontinuities in the Earth's interior. Discuss with diagram the passages of P and S waves inside the Earth. \pts{14}

\bigskip\noindent{\large\bfseries SECTION --- B}
\rule{\linewidth}{0.4pt}\smallskip

\noindent\textbf{Question 5.} Write short notes on the following: \pts{$2 \times 7 = 14$}
\begin{parts}
  \item Discordant igneous structures.
  \item Strike and dip and right-hand rule of recording of strike.
\end{parts}

\medskip\noindent\textbf{Question 6.} Discuss the following: \pts{$2 \times 7 = 14$}
\begin{parts}
  \item Outlier and inlier with neat diagram.
  \item Types of sedimentary rocks.
\end{parts}

\bigskip\noindent{\large\bfseries SECTION --- C}
\rule{\linewidth}{0.4pt}\smallskip

\noindent\textbf{Question 7.} Briefly discuss the following: \pts{$4 \times 3\frac{1}{2} = 14$}
\begin{parts}
  \item Strike and dip joints.
  \item Basic features of contour lines.
  \item Origin of Earth's second atmosphere.
  \item Causes of earthquake.
\end{parts}

\medskip\noindent\textbf{Question 8.} Write about each of the following in 2 to 3 sentences. Give suitable diagrams wherever necessary: \pts{$14 \times 1 = 14$}
\begin{parts}
  \item Contour map of conical hill
  \item Inflection point
  \item Axial plane
  \item Antiformal syncline
  \item Overturned antiform
  \item Outcrop and incrop
  \item Relationship between horizontal bed and contour lines
  \item The biggest mass extinction event in the Earth's history was during the \underline{\hspace{2cm}}.
  \item Nebular hypothesis was proposed by \underline{\hspace{2cm}}.
  \item Stratovolcano
  \item Index fossil with a suitable example
  \item Name two metamorphic rocks
  \item Half life
  \item Earth's magnetic field is generated by \underline{\hspace{2cm}}.
\end{parts}

\vfill\rule{\linewidth}{0.4pt}
\begin{center}\small
\href{https://bhu-exam.vercel.app/}{Click here to visit our website}\\[4pt]
\href{https://me-aryan.vercel.app/}{Click here to visit our contributor}\\[4pt]
\href{https://quashan-me.vercel.app/}{Click here to visit our contributor}
\end{center}
\end{document}
